\section{指针转换：拆除前的准备}

\begin{warnbox}
在调用 \texttt{vmm\_unmap\_identity()} 之前，必须将所有仍指向低地址的指针
转换为 high-half 地址。否则拆除 identity mapping 后，解引用这些指针会
触发 Page Fault。
\end{warnbox}

在我们的内核中，需要转换的关键指针包括：

\subsection{\texttt{g\_mb2\_info}：Multiboot2 信息结构}

GRUB 传入的 Multiboot2 info 指针是物理地址（由 \texttt{boot.asm} 从 \reg{EBX} 保存）。
在 \texttt{vmm\_init()} 建立 high-half mapping 之后，需要将其转换：

\codefile{src/kernel/core/kmain.c}
\begin{lstlisting}[style=cstyle, caption={kmain.c：转换 Multiboot2 指针}]
/* vmm_init 之后、vmm_unmap_identity 之前 */
g_mb2_info = (const void*)PHYS_TO_VIRT(
    (uint32_t)(uintptr_t)g_mb2_info);
\end{lstlisting}

如果忘记转换会怎样？
\begin{itemize}
  \item \texttt{g\_mb2\_info} 仍指向 $\sim$\hex{00010000}（物理地址）
  \item 此时 identity mapping 还在，暂时不会出错
  \item 但 \texttt{vmm\_unmap\_identity()} 之后，PD[0] 被清除
  \item 下次 shell 命令访问 \texttt{g\_mb2\_info} → \hex{00010000} → Page Fault!
\end{itemize}

\subsection{PMM bitmap 指针}

PMM bitmap 后端在 \texttt{pmm\_init()} 时将 bitmap 放置在物理内存中。
由于 bitmap 后端内部存储的是虚拟地址指针（通过 identity mapping 期间可以等价于物理地址），
在 \texttt{vmm\_init()} 建立 high-half mapping 后，PMM 内部的 bitmap 指针
需要通过直接映射区的 high-half 地址来访问。

\begin{thinkbox}
如果 PMM bitmap 存储在物理地址 \hex{00200000}，在 identity mapping 拆除后，
内核应该通过哪个虚拟地址访问它？

答：\hex{C0200000}（$= \hex{00200000} + \hex{C0000000}$）
\end{thinkbox}

% ============================================================================

